\begin{WideTable}
    \caption{Inference test categories that exercise behavioral predicates.}
    \label{tab:inference-tests}
    \centering
    \small
    \begin{tabularx}{\linewidth}{l X}
      \hline
      \textbf{Category} & \textbf{Representative cases} \\
      \hline
      Direct calls &
        Single and chained calls to a spec'd helper; inferred specs use
        \texttt{ensures\_of} and \texttt{aborts\_of} to abstract over the
        callee's effect rather than inlining its body. \\
      \hline
      Mutual recursion &
        Mutually recursive \texttt{is\_even}/\texttt{is\_odd}; inference
        breaks the cycle by phrasing each spec in terms of the other's
        behavioral predicate. \\
      \hline
      Higher-order parameters &
        A function that takes a closure and applies it; the inferred spec
        abstracts the closure via \texttt{ensures\_of}/\texttt{aborts\_of}
        on the parameter, with no commitment to a particular implementation. \\
      \hline
      Stored-closure dispatch (enum) &
        Calculator with an enum variant carrying a closure field;
        dispatch through the variant is summarized by behavioral
        predicates on that field. \\
      \hline
      Stored-closure dispatch (struct field) &
        Vault with a \texttt{Strategy} closure stored in a struct field on
        Aptos framework types; inference emits \texttt{modifies\_of} for
        the field together with the field-form behavioral predicates. \\
      \hline
      Pre/post state labels &
        \texttt{move\_to}, \texttt{move\_from}, mutable resource indexing;
        inferred specs distinguish \texttt{@pre} from post-state on the
        memory the predicate observes. \\
      \hline
      Intermediate state labels &
        Sequenced state-modifying calls (e.g.\ remove-then-publish);
        inferred specs introduce existentially quantified intermediate
        labels and pin them to defining behavioral predicates. \\
      \hline
    \end{tabularx}
\end{WideTable}


\Section{Specification Inference}
\label{sec:inference}

\SubSection{How it Works}
\label{sec:validation-overview}

\MVP's \emph{specification inference} \cite{GrieskampEtAl26Inf} based on weakest-precondition (\WP) analysis~\cite{dijkstra1975wp,BL05} is made feasible and modular via behavioral predicates.
\WP inference takes arbitrary Move code and, for each function, computes conditions that are both necessary and sufficient for the implementation's behavior.
Spec inference uses behavioral predicates to express the effect of calling out to another function, whether through a function value or a direct call.
It reverses the translation seen for $invoke_\tau$ in Sec.~\ref{sec:enc-invoke}, introducing the behavioral predicates into the \WP to express the call. Without this modular abstraction, inferred specs would be impractically large.

A full description of \MVP's \WP analysis is out of scope of this paper. What is important is that, besides the notorious loop invariant problem for \WP, specification inference for Move is precise and complete.
This serves two purposes for this work, demonstrated in the following subsections: reducing the specification burden for compliance checking of function values (Sec.~\ref{sec:validation-amm}), and providing a validation pipeline that confirms the encoding is internally consistent (Sec.~\ref{sec:validation-tests}).

\SubSection{Compliance Checking}
\label{sec:validation-amm}

Recall from Sec.~\ref{sec:amm} that the !Pool! struct carries a !pricing! closure constrained by a no-abort invariant: the pricing function must not abort for any inputs in any state.
Two concrete pricing implementations were shown: !product!, which is compliant (it never aborts), and !with_fee!, which is not (it aborts when the owner's !Fee! resource is absent).
Spec inference makes this compliance distinction machine-readable without any manual annotation.

Given the implementation of !with_fee!, \WP inference easily derives the aborts conditions (we leave out the ensures, which it also determines):

% @formatter:off
\begin{MoveBox}
  spec with_fee {
    aborts_if [inferred] !exists<Fee>(owner);
    aborts_if [inferred] Fee[owner].bps > 10000;
    ...
  }
\end{MoveBox}
% @formatter:on

Similarly, for the lambda example in Sec.~\ref{sec:find}, we were required to attach a spec to a trivial lambda expression of the form !|x| x == 0!.
Obviously, inference can help to avoid this.
We are currently investigating relaxing the requirement of spec blocks for compliance-checked functions in those cases where they can be reliably inferred (no recursion and no loops); however, this is orthogonal to the work in this paper.


\SubSection{Validation Pipeline and Test Coverage}
\label{sec:validation-tests}

Inferred conditions are labeled \texttt{[inferred]} and feed directly back into \MVP for verification, creating a pipeline: code $\leadsto$ code+specs $\leadsto$ VC, with verification expected to succeed.
This pipeline acts as a \emph{differential-testing oracle} for the extension described in this paper: inference and verification are independent implementations which meet only at the semantics of behavioral predicates and state labels, so a failure indicates a defect in one of them --- conditions inferred unsoundly or not tightly, or axioms too weak to discharge them.
Tab.~\ref{tab:inference-tests} maps the categories of the inference test suite~\cite{InferenceTestsuite} --- currently 27 cases with over 550 inferred conditions --- to the HOF patterns introduced in this paper.
All categories pass end-to-end, evidence that the encoding is internally consistent and that the axioms are strong enough to discharge \emph{tight} specifications: usable, not merely sound.
Mutual recursion deserves a note: inference breaks the specification cycle by phrasing each function's spec in terms of the other's behavioral predicate, which verification resolves against the axioms of Sec.~\ref{sec:enc-bp}.

\SubSection{Empirical Status}
\label{sec:empirical}

The extension described in this paper is implemented in the production \MVP shipped with the Aptos toolchain; the paper describes Move language version 2.4 and \MVP as of the aptos-core commit pinned in~\cite{PaperExamples, InferenceTestsuite}.
The closure portion of the prover's test suite comprises 19 test modules with about 240 verified functions and 66 expected-error diagnoses, including the examples of this paper; together with the inference suite, these run in continuous integration on every change to the prover.
The tests are required to be stable across changes to the prover and underlying tools like Z3, and complete in under 5 minutes on standard GitHub runners, with a default timeout of 40 seconds per verification condition.
A benchmark over a larger corpus of higher-order Move contracts does not yet exist: higher-order Move shipped in late 2025, and a corpus is only now forming, partly in closed-source projects.




%%% Local Variables:
%%% mode: latex
%%% TeX-master: "main"
%%% End:
